%\documentclass{article}
%\begin{document}
  \huge{\bfseries{PART A}}
  \large 
  \section{Problem Identification}
    \subsection{Problem statement}
      Attention Deficit Hyperactivity Disorder (ADHD) prevalence is increasing
      and the sufferers of this condition struggle with everyday school and
      work due to having difficulties with focus and attention retention.

    \subsection{Key issues}
      There are many behavioural disorders that affect human comfort and living
      standards. One of the most common behavioural disorders is 
      Attention Deficit Hyperactivity Disorder (ADHD). ADHD is characterised by 
      the decline in focus over periods of time due to low  brain activity
      prompting the brain to become hyperactive to reach healthy/ activity 
      levels \cite{timothy}. Sufferers require external stimulation 
      supplementary to the low stimulus activity to maintain focus.
      Performing motor activity has been shown as a reliable way
      to reduce hyperactivity and distractibility in ADHD patients 
      \cite{azrin}).\\ \\
      Current testimony indicates that the primary issue is that ADHD sufferers
      tend to lose attention when presented with information which can have
      negative effects on school performance as well as workplace relations
      as distractible workers are regarded as poor performers.

    \subsection{Affected population}
      ADHD has not shown any bearing on geographical location, gender or
      ethnicity. ADHD is correlated with age showing that symptoms tend to
      decrease with increasing age of patients \cite{josephine}. The
      symptoms are primarily present in groups that are required to stay
      stationary and give attention to one activity for an extended period
      of time. \\ \\
      By this it is concluded that the affected population is children in 
      school and young adults working white collar jobs.

    \subsection{Short commings of existing solutions}
      Existing solutions include fidget toys which altough effectiveness is
      somewhat proven \cite{amanda}. These toys require the use of 1 or 2 
      hands, which can be a hindrance when considering the use of a sufferer's
      hands in a classroom or work setting. In the class room a fidget toy
      taking up a hand can make it difficult for the student to take notes, or
      a worker to use tools or type on a keyboard. \\ \\
      Another concern is that of noise, many existing fidget toys can produce
      noise especially when used vigorously as is common with children.
      These toys also have a tendency to degrade over time causing the device to
      produce even more noise. The presence of noise makes these devices
      impractical for use in a classroom as well as meetings or any other noise
      sensitive environments where sufferers will be required to focus.


    \subsection{Evidence to support the investment decision}
      According to a 2016 study on factors that promote resilience in youth with
      ADHD symptoms \cite{melissa}. The single bigest factor is social
      acceptance. Thus, the priority should be to symptomatically treat ADHD
      symptoms while minimizing the effect on the surrounding social
      environment. \\ \\
      A 2011 study done on the efficacy on fidget toys found that children who
      suffer from ADHD that were given sensory challenging tasks that were
      measured on how much of the time spent solving the problem was spent on
      tasks and how much time was spent doing off-task activities. At the lowest
      end when the children were told to use a fidget toy they spent 7.9 \% less
      time doing off-task activities and on the high end time spent doing
      off-task activities decreased by 52.8\% \cite{amanda}.  \\ \\
      From this study it can be confidently stated that fidget toys do indeed
      work. Following this another study done in 2018 shows that the use of
      alternative seating arrangements in classrooms  \cite{julaila}
      . The study tested the use of 5 alternative seating arrangements,
      namely Therapy Balls, Disc 'O' Sit cushions, Therapy Balls with back rest,
      and Stability balls. All 4 alternative seating methods showed significant
      increases in the test groups on task behaviours, seated practice as well
      as well as increased comfort \cite{julaila}. \\ \\
      from the above discussed studies it is abundantly clear that this is a
      necessary product to increase the comfort of human living as well as
      helping the rapidly increasing population of ADHD sufferers \cite{elie}.
       

    \subsection{Reframing the problem}
    There is a gap in the market for ADHD fidget devices. Current solutions are
    either fidget toys or stationary standing desks or therapy balls. Although
    these devices are proven to be helpful they all have either noise issues,
    require you to sacrifice the use of one or both hands or is meant to be used
    stationarily meaning the device cannot be carried around restricting
    portability making them inadequate for busy school, or young professional
    life. A product needs to be designed that allows adequate stimulation for
    ADHD sufferers, while having both hands free that is easily portable between
    classrooms and boardrooms.



  \section{Objectives and Operational requirements}
    \subsection{business system's requirements} 
      \tikzset{parent/.style={align=center,text width=10cm,fill=green!20,rounded corners=2pt},
      child/.style={align=center,text width=3.5cm,fill=green!50,rounded corners=6pt},
      grandchild/.style={fill=pink!50,text width=2.1cm} 
      }
      \begin{figure}[H]
      \begin{forest}
        for tree={%
          thick,
          drop shadow,
          l sep=0.6cm,
          s sep=0.8cm,
          node options={draw,font=\sffamily},
          edge={semithick,-Latex},
          where level=0{parent}{},
          where level=1{
            minimum height=1cm,
            child,
            parent anchor=south west,
            tier=p,
            l sep=0.25cm,
            for descendants={%
              grandchild,
              minimum height=0.6cm,
              anchor=150,
              edge path={
                \noexpand\path[\forestoption{edge}]
                (!to tier=p.parent anchor) |-(.child anchor)\forestoption{edge label};
              },
            }
          }{},
        }
        [To Provide comfort and stimulation of ADHD sufferers
          [Affordability
            [Affordable\\Materials
              [Affordable\\Manufac-\\turing\\Process]
            ]
          ]     
          [Portability
            [Lightweight\\Materials
              [Compact\\Design
                [Easily\\Transportable
                  [Easy\\To\\Setup]
                ]
              ]
            ]
          ]
          [Durability
            [Durable\\Materials
              [Easily\\Repairable\\Design]
            ]
          ]
          [{Professional\\Design}
            [Singular\\Professionally\\Acceptable\\Color
              [Purely\\Functional\\Form factor]
            ]
          ]
        ]
      \end{forest}
        \caption{Operational Objectives Outline}
      \end{figure}
      %\begin{center}
      %  \textit{Figure 1 : Operational Objectives Outline} 
      %\end{center}
      The above diagram shows an outline of the required objectives for the
      business to complete and what the main concern for achieving that 
      objective should be. These points will be discussed further in later
      parts of this document. \\ \\
      The Operational drivers that will be required : 
      \begin{enumerate}
        \item Pharmacological dealers \\
          This product will be difficult to sell from general dealers thus the
          best outlet should sell pharmacological good, examples of this is
          Clicks and Dischem.
          The health and distribution of these dealers are
          necessary to the success of this business. 
        \item Online dealers \\
          As the modern online dealerships such as Takealot and Temu rapidly
          increase in popularity. The health and logistics capabilities of these
          businesses will drastically aid in the logistics system of the
          business. Takealot and Temu (Online Dealers) will be approached as
          po
        \item Aluminum Suppliers \\ 
          The best material to fit the laid out objectives for the product
          frame is aluminum. The health of the aluminum industry is a required
          operational driver as the production of aluminum would be out of scope
          to handle for this product's production process.
        \item Nylon Suppliers \\
          The best durable material for the sitting air bag as well as the
          attached carry bag is nylon. Thus the health of the nylon industry is
          a required operational driver of the products production process.
        \item Silicone Spray suppliers \\
          To properly seal the nylon inflatable bag the nylon fabric must
          be coated with an air impermeable material, This silicone spray is a
          required operational driver for the production process of this
          product.
        \item Spray Paint suppliers \\
          To achieve the professional form-factor, as well as to increase
          product durability Paint is required to coat the product. Thus The
          health of the Spray Paint industry is an operational driver of this
          product's production process.
        \item Customer Movement \\
          One of the product's market opportunities is the ability to be
          portable. Any global event that will limit the mobility of
          customers such as a global pandemic or mass shift to working from home
          will be a threat to the success of the product as this type of
          situation will reduce the market gap in portable chair attachments.
        \item ADHD Prevalence \\
          This system requires that ADHD is a problem that customers experience
          there is currently no known cure or full proof symptomatic
          treatment for ADHD but any such medical advances will require a
          re-evaluation of the system.
      \end{enumerate} 

    \subsection{required capabilities}
      \begin{enumerate}
      \item Material Sourcing \\
        A required capability of the system is the ability to source adequate
        materials for the production process of the product. \\  
        Having proper material sourcing capabilities addresses the operational
        objectives of \textbf{affordable materials}, 
        \textbf{lightweight materials} and \textbf{durable materials}. \\
          The system requires the ability to source Aluminum and textile
          nylon. Both of these materials are relatively cheap in their respective
          industries as well as known for being lightweight. The durability of
          nylon is notorious among textile producers. Aluminum provides adequate
          durability for this application as it is often used in the production
          of bicycle frames. The physical properties of these materials addresses
          \textbf{all material objectives} as well as the \textbf{easily
          transportable} objective. \\
        \begin{figure}[H]
        \begin{forest}
          for descendants={anchor=east,child anchor=west,grow=east,
            for descendants={anchor=east, child anchor=west,grow=east,},
          }, 
          grow=east,anchor=west,parent anchor=east,
          for tree={fill=green!20},where={n()<=3}{draw=black}{},
          where level=1{fill=green!50,}{},
          where level=2{fill=pink!50,}{},
          l sep=1cm,
          s sep=0.8cm,
          edge={semithick,-Latex},
          drop shadow,
          %
          [{Material Sourcing addresses},
            [{Affordability},
              [{Affordable Materials}]
            ]
            [{Portability},
              [{Lightweight Materials},]
              [{Easily Transportable},]
            ]
            [{Durability},
              [{Durable Materials},]
            ]  
          ] 
        \end{forest}
          \caption{Material Sourcing addresses}
      \end{figure}
        %\begin{center}
        %  \textit{Figure 2: Material Sourcing addresses}
        %\end{center}
        The above diagram show how the material sourcing capability addresses the goals
          and objectives laid out by the system objectives.

      \item Optimized Production \\
        This business system requires the capability to effectively produce the
        necessary components of the product. \\
        Having a optimized production system will address and aid to achieve the
        objective of \textbf{affordable manufacturing}\\
        \begin{figure}[H]
        \begin{forest}
          for descendants={anchor=east,child anchor=west,grow=east,
            for descendants={anchor=east, child anchor=west,grow=east,},
          }, 
          grow=east,anchor=west,parent anchor=east,
          for tree={fill=green!20},where={n()<=3}{draw=black}{},
          where level=1{fill=green!50,}{},
          where level=2{fill=pink!50,}{},
          l sep=1cm,
          s sep=0.8cm,
          edge={semithick,-Latex},
          drop shadow,
          %
          [{Optimized Production addresses},
            [{Affordability},
              [{Affordable Manufacturing}]
            ]
          ] 
        \end{forest}
          \caption{Optimized Production addresses}
        \end{figure}
        %\begin{center}
        %  \textit{Figure 3: Optimized Production addresses}
        %\end{center}
        The above diagram show how the optimized production capability addresses the goals
          and objectives layed out by the system objectives.
      
      \item Simple and standardized parts and interfaces. \;
        Designing the product with simple and standardized interfaces with
        skew-in capabilities\\
        Having simple and standardized parts and interfaces will be crucial to
        achieve the objectives of \textbf{easily repairable} and
        \textbf{affordable manufacturing}. Using skew-in interfaces will aid with
          the goal \textbf{Durability} due to skew-in interfaces having better
          force dissipation as as well as skew-in interfaces being cheap to
          manufacture achieving the objective of 
          \textbf{affordable manufacturing}. \\
          By utilizing standard parts with no extra non-essential features aid
          in the objective of appearing professionally acceptable ny having a
          \textbf{Purely Functional form factor}. \\
        \begin{figure}[H] 
        \begin{forest}
          for descendants={anchor=east,child anchor=west,grow=east,
            for descendants={anchor=east, child anchor=west,grow=east,},
          }, 
          grow=east,anchor=west,parent anchor=east,
          for tree={fill=green!20},where={n()<=3}{draw=black}{},
          where level=1{fill=green!50,}{},
          where level=2{fill=pink!50,}{},
          l sep=1cm,
          s sep=0.8cm,
          edge={semithick,-Latex},
          drop shadow,
          %
          [{Simple, Standardized parts},
            [{Affordability},
              [{Affordable Manufacturing}]
            ]
            [{Durability}
              [{Easily repairable}]
            ]
            [{Professional Design}
              [{Purely Functional Form factor}]
            ]
          ]
        \end{forest}
        \caption{Simple Standardized parts and interfaces addresses}
        \end{figure}
         %\begin{center}
         %  \textit{Figure 4: Simple Standardized parts and interfaces addresses}
         %\end{center}
        The above diagram show how the use of standardized parts capability addresses the goals
          and objectives layed out by the system objectives.
      
      
      \item Simple interlocking Joints \\
        The business system requires the production of simple interlocking
          joints that utilize standardized screw-in interfaces. \\
          The addition of simple interlocking joints will achieve the objectives
          of having a \textbf{compact design} as parts can be folded to make the
          product more compact and more \textbf{easily transportable}. Due to
          the joints being interlocking it also makes to product \textbf{easy to
          setup} \\
        \begin{figure}[H]
        \begin{forest}
          for descendants={anchor=east,child anchor=west,grow=east,
            for descendants={anchor=east, child anchor=west,grow=east,},
          }, 
          grow=east,anchor=west,parent anchor=east,
          for tree={fill=green!20},where={n()<=3}{draw=black}{},
          where level=1{fill=green!50,}{},
          where level=2{fill=pink!50,}{},
          l sep=1cm,
          s sep=0.8cm,
          edge={semithick,-Latex},
          drop shadow,
          %
          [{Interlocking Joints addresses},
            [{Portability},
              [{compact design}]
              [{Easily transportable}]
              [{Easy to setup}]
            ]
          ]
        \end{forest}
          \caption{Simple Interlocking Joints addresses}
        \end{figure}
         %\begin{center}
         %  \textit{Figure 5: Simple Interlocking Joints addresses}
         %\end{center}
        The above diagram show how the use interlocking joints capability addresses the goals
          and objectives layed out by the system objectives.

      \item Painting \\
        The project require that the system has capabilities to paint the
          product professionally appealing colours. As well as to be able to apply
          an air impermeable silicone coating.\\
          Being able to paint the device a flat colour will aid to achieve the
          \textbf{Professional\\Design} goal. \\
        \begin{figure}[H]
        \begin{forest}
          for descendants={anchor=east,child anchor=west,grow=east,
            for descendants={anchor=east, child anchor=west,grow=east,},
          }, 
          grow=east,anchor=west,parent anchor=east,
          for tree={fill=green!20},where={n()<=3}{draw=black}{},
          where level=1{fill=green!50,}{},
          where level=2{fill=pink!50,}{},
          l sep=1cm,
          s sep=0.8cm,
          edge={semithick,-Latex},
          drop shadow,
          %
          [{Painting},
            [{Professional Design},
              [{Singular Professionally acceptable colour}]
            ]
          ]
        \end{forest}
          \caption{Painting addresses}
        \end{figure}
        %\begin{center}
        %   \textit{Figure 6: Painting}
        %\end{center}
        The above diagram show how the use Painting capability addresses the goals
          and objectives layed out by the system objectives.
      \end{enumerate}
    \subsection{Physical embodiments for the functions}
      \begin{enumerate} 
        \item Material sourcing physical embodiments \\
          The Material sourcing capability of the system requires the physical
          embodiments of a reliable logistics system for the transport
          of raw material to a centralized factory. Thus Transport hardware is
          required, requiring trucks and drivers. \\
        \item Optimized Production system physical embodiments \\
          For an optimized manufacturing process it is required that automation
          through industrial machinery is maximized while minimizing
          hand-crafting of parts. For this to be effective the shaping of
          aluminum parts will have to be done with machinery thus requiring the
          physical embodiment of a milling machine, drilling machine as well as
          specialized drill-bits for tapping. There will also be a requirement
          for mend different parts to one another permanently for this a welding
          machine is required. Due to the difficulty of welding aluminum a
          TIG-welder will be required. For the manufacturing to be as
          cost-effective as possible software should be used instead of a human
          operator for the TIG-welder. The required software is Miller Welds
          provided by the company of the same name. \\ For the production of the
          attached transport bags the required nylon will be transported with
          the same method as the aluminum. Nylon fabric can be bought as
          sheets thus only requiring cutting and weaving. For this there will be
          a requirement for a sewing machine as well as a fabric cutter. To
          effectively and quickly inflate the internal air bags a relatively
          small air compressor will be required. 
        \item Standardized parts and interfaces \\
          All requirements for standardizing parts will already have been
          covered in the Optimized Production system section.
        \item Simple interlocking joints \\
          The process of creating an interlocking joint requires the use of
          springs. To self-produce the required spring is unnecessarily
          expensive thus sourcing spring on their own will be done. This
          physical embodiment is already provided for by the material sourcing
          section.
        \item Painting \\
          For the painting process only a single black colour will be used. The
          paint needs to be sourced using the same logistics system as proposed
          by the material sourcing part of the physical embodiments. To
          cost-effectively paint large quantities of parts with the same colour
          will require spray painting tools. These tools consist of both
          consumable masks for operators as well as the \"spray gun\".
        \item Silicone coating \\
          Silicone coating needs to be sourced to be able to seal the internal
          air-bag, Silicone coating can be applied using a \"Spray gun\" similar
          to that which is used in the Painting process.
        \item Air Inflation Valve's  \\
          To allow the air-bag to be re-pressurised after significant use, a simple
          Air Inflation valve is required.
      \end{enumerate}
    \subsection{Needs of the business system}
      As outlined by the above-mentioned parts the business system requires many
      different systems to work together. 
      \begin{enumerate}
        \item Measure of Effectiveness (MOE) \\
          To measure the effectiveness of the system as a whole. The measure of
          effectiveness needs to directly address the core problem at hand as
          well as address the success of the business. Therefore the measure of
          effectiveness should be measured by the percentage increase in
          observed on task behaviour of individuals using the device versus
          their baseline on-task behaviour, averaged across all units sold
          multiplied by number of units sold.
        \item Measure of Performance (MOP) \\
          To measure the effectiveness of the subsystems the following measure
          will be taken.
          \begin{enumerate}
            \item \textbf{Average variance from the target material cost. }
              the lower the variance is, the more reliably the materials can be
              sourced within the allocated budget. Directly addressing the
              "Affordable Materials" objective. 
            \item \textbf{Frequency and fraction of On-time material
              deliveries.} The higher the frequency and fraction of on-time
              deliveries directly measures the effectiveness of the logistics
              system. 
            \item \textbf{Material quality conformance fraction.} By measuring
              the fraction of materials that conform to the quality standard
              required will measure the efficiency and trust of the material
              sourcing partners.
            \item \textbf{Paint time and Paint consumption.} By measuring the
              time and volume paint required to fully paint one unit of the
              product to minimize both, directly addressing the "Affordability"
              goal.
            \item \textbf{Per-Part manufacturing time and material cost.} By
              measuring the manufacturing time and material cost per part of the
              device possible bottlenecks could be eliminated as well as
              individual parts simplified. Also analyzing manufacturing time to
              simplify and streamline the manufacturing process. Directly
              addressing the "Affordable manufacturing" objective.
          \end{enumerate}
        \item Technical Performance Measure (TPM) \\
          \begin{enumerate}
            \item \textbf{Number of assembly steps.} Since the device will have
              to be frequently transported the number of steps required to
              assemble the product will directly address the "Standardized parts"
              and "Simple Interlocking joints" physical embodiments. The TPM
              will be to have a maximum number of allowed assembly steps.
            \item \textbf{Folded Dimensions and Weight.} Measuring the folded
              dimensions as well as the weight of the device with the aim to get
              the device as small and light as possible. This Directly aids to
              achieve the "Portable" goal. The TPM will be to have a target
              weight and dimension that the product has to conform to.
            \item \textbf{Joint duty-cycle and weight limit.} With the product's goal of 
              durability, the duty-cycle of the joints as well as their weight
              limits. The TPM will be to have a target duty-cycle number and a
              weight limit for each joint.
            \item \textbf{Paint Adhesion strength.} Due to the goal of
              professional design, the paint and appearance of the product
              needs to be maintained over time requiring that the paint does not chip
              off over time. The TPM will be to have a minimum score achieved on
              the Standardized paint adhesion test.
            \item \textbf{Silicone coating sealing quality.} Due to the required
              internal air-bag for better ADHD stimulation the performance of the
              air-bag is measured by how well the internal silicone coating seals
              the bag so that it does not leak air. This is determined by
              measuring the amount of time it takes for the bag to deflate when
              under a specified amount of pressure.
          \end{enumerate}
      \end{enumerate}
 %\end{document}
